% !PRINTED 607
\weberpage{607}

% [German fragment carries a correction note: the running head of printed 607
% ("§. 179. | Reelle Radicale. | 607") had been transcribed as a section heading and was
% removed; running heads are excluded page furniture. Mirrored here for symmetry.]
\quad We assume a \gesperrt{real domain of rationality} $\Omega$,
and take in it a normal equation $g(t) = 0$,
i.\,e. an irreducible equation whose roots are all rationally
expressible in terms of an arbitrary one of them, $\varrho$. If one
of the roots is real, then all the others must be real
as well, and $g(t)$ thus has either only real or only pairwise
conjugate imaginary roots. $g(t) = 0$ can be the \gesperrt{Galois}
resolvent of any given equation $F(x) = 0$, whether irreducible
or reducible, and if therefore $g(t)$ has only
real roots, then $F(x)$ too can have only real roots, % [sic: the German page prints "relle" for "reelle"]
because the roots of $F(x)$ are rationally representable
in terms of a root of $g(t)$. If among the roots of $F(x)$ there are
imaginary ones, then $g(t)$ has only imaginary roots. But if $F(x)$
has only real roots, then the roots of $g(t)$ are also real,
because they can indeed be expressed rationally in terms of the
roots of $F(x)$.

\quad Now if $g(t)$ is reduced by the adjunction of a root $\varepsilon$ of an
irreducible equation $\chi = 0$ \gesperrt{whose degree is a prime},
then, as we have seen in \S.\,157, all the
roots of $\chi$ are contained in $\Omega(\varrho)$, and if therefore
$\varrho$ is real, then all the roots of $\chi$ are real.

\quad 1. \gesperrt{A normal equation with real roots can
therefore be reduced only by such irreducible
equations of prime degree as have only
real roots.}

\quad We now ask whether a reduction of the normal equation
can be effected by the adjunction of a real radical $\sqrt[p]{a}$.
We may here assume that the degree $p$ of the radical
is a prime, because every radical can be reduced to a series of
extractions of roots of prime degree. We can also
presuppose that $a$ is not the $p$th power of a rational
quantity, because otherwise the real root $\sqrt[p]{a}$ would be rational.

\quad Under this assumption, as we now wish to prove, % [sic: the German page prints "Vorraussetzung" for "Voraussetzung"]
$x^p - a$ is irreducible. For if we denote one, e.\,g. the real
value $\sqrt[p]{a}$, by $\alpha$, and an imaginary $p$th root of unity by $\varepsilon$,
then
\weq{1}{\alpha,\ \varepsilon\alpha,\ \varepsilon^2\alpha \ldots \varepsilon^{p-1}\alpha}
%CONT>>
